\documentclass[10pt,a4paper]{amsart}
\usepackage[margin=19mm]{geometry}
\usepackage{amsmath,amssymb,mathtools}
\usepackage[scr=rsfs,cal=euler]{mathalfa}
\usepackage{xfrac}
\usepackage[unicode,hidelinks,bookmarks=false]{hyperref}
\newcommand{\ie}{i.e.\ }
\newcommand{\cf}{cf.\ }
\newcommand{\resp}{respectively\ }
\newcommand{\Subsection}{\subsection}
\providecommand{\nobreakdash}{\nobreak\hbox{-}}
\setlength{\emergencystretch}{2em}
\multlinegap0pt
\usepackage{amssymb}
\usepackage{xcolor}
\usepackage{tikz}
\usepackage[shortlabels]{enumitem}
\usepackage[normalem]{ulem}
\newtheorem{proposition}{Proposition}
\newcommand{\ZZ}{\mathbb{Z}}
\newcommand{\one}{\mathbb{1}}
\newcommand{\zero}{\mathbb{0}}
\title{Note on adjacency spectra for twice punctured hypercubes}
\author{Jang Soo Kim, Victor Reiner, Laura Schoenzeit}
\date{Source: arXiv:2608.27887v1 (CC BY 4.0)}
\begin{document}
\maketitle

\noindent\emph{Source and licence.} Note on adjacency spectra for twice punctured hypercubes. Version arXiv:2608.27887v1 (CC BY 4.0), distributed by arXiv under the Creative Commons Attribution 4.0 International licence, http://creativecommons.org/licenses/by/4.0/, as recorded in the arXiv licence metadata for this exact version and shown on the abstract page https://arxiv.org/abs/2608.27887v1. Full licence notice: \texttt{licenses/arxiv-2608.27887-CC-BY-4.0.txt}.\par\smallskip

\noindent\textit{Editorial scope.} Counted unit: the article's proposition prop:difference-eigenvectors (source0.tex lines 381--390). The article's complete original proof of it is delivered with it (source0.tex lines 391--439, one \texttt{proof} environment of 49 lines). No mathematics is written, repaired or re-derived: the statement and the proof are the article's own lines, in the article's own notation, and no retained line is substituted or re-flowed.  Retained uncounted as the source-local context the proof needs: lines 335--338.  Nothing was omitted from any retained span, so the package carries no omission record.  No retained line contains a citation, so the wrapper carries no bibliography and no bibliography entry.  No retained line contains a figure environment, an image-inclusion command or drawing code, so no image is delivered and the package declares no asset dependency.  Editorial formatting only, in addition: an AMS \texttt{amsart} preamble that restates the article's own declarations and macros used by the retained lines, namely the environments \texttt{proposition} on separate counters, together with the standard packages those lines need.  The retained environments print in this document as Proposition 1 in order of appearance, whereas the article prints the same items with the numbering of its own full document.  The complete native source of the article as distributed in the arXiv e-print for this exact version is bundled unchanged as \texttt{sources/208747/source0.tex}.
\begin{equation}
\label{eqn:eigenvector-equation-for-Eu}
A_{Q_n}(E_u)=(n-2k) E_u.
\end{equation}

\begin{proposition}
\label{prop:difference-eigenvectors}
Let $u,v$ be in $\ZZ_2^n$, 
with $u \neq v$ and $\omega(u)=\omega(v)=k$ for $1 \leq k \leq n-1$.
Then their difference vector $E_u-E_v$
lies in $\hat{\hat{V}}$,
and gives a common $(n-2k)$-eigenvector
for all three matrices $A_{\hat{\hat{Q}}_n}, A_{\hat{Q}_n}, A_{Q_n}$.

\end{proposition}

\begin{proof}
The assertion that the difference vector $E_u-E_v$
lies in $\hat{\hat{V}}$
follows from these equalities:
\begin{align*}
(E_u)_{\zero_n}&=(-1)^{u \cdot \zero_n}=(-1)^0=(-1)^{v \cdot \zero_n}=(E_v)_{\zero_n},\\
(E_u)_{\one_n}&=(-1)^{u \cdot \one_n}=(-1)^k=(-1)^{v \cdot \one_n} =(E_v)_{\one_n}.
\end{align*}
To prove the common eigenvector assertion, first note that given 
any $w$ in
$\ZZ_2^n$, one has
$$
(A_{Q_n} E_u)_w
=\sum_{\substack{w' \in \ZZ_2^n:\\
d(w,w')=1}} (E_u)_{w'}
=\sum_{\substack{w' \in \ZZ_2^n:\\
d(w,w')=1}} (-1)^{u \cdot w'}.
$$
Similarly, for \( w\in \ZZ_2^n \setminus \{\zero_n\} \), we have
\begin{align*}
(A_{\hat{Q}_n} (E_u-E_v))_w
&=\sum_{\substack{w' \in \ZZ_2^n \setminus \{\zero_n\}:\\
d(w,w')=1}}
[(-1)^{u \cdot w'}
- (-1)^{v \cdot w'}]
\overset{(a)}{=}\sum_{\substack{w' \in \ZZ_2^n:\\
d(w,w')=1}} 
[(-1)^{u \cdot w'}
- (-1)^{v \cdot w'}]\\
&=(A_{Q_n} (E_u-E_v))_w
\overset{(b)}{=}(n-2k) (E_u-E_v)_w,
\end{align*}
where equality (a) used the fact that $(-1)^{u \cdot \zero_n}=(-1)^{v \cdot \zero_n}$, and (b) used the eigenvector equation \eqref{eqn:eigenvector-equation-for-Eu}.  
Similarly, since $(-1)^{u \cdot \one_n}=(-1)^k=(-1)^{v \cdot \one_n}$, 
for any $w\in \ZZ_2^n \setminus \{\zero_n,\one_n\}$, we have
\begin{align*}
(A_{\hat{\hat{Q}}_n} (E_u-E_v))_w
&=\sum_{\substack{w' \in \ZZ_2^n \setminus \{\zero_n,\one_n\}:\\
d(w,w')=1}} 
[(-1)^{u \cdot w'}
- (-1)^{v \cdot w'}]
=\sum_{\substack{w' \in \ZZ_2^n:\\
d(w,w')=1}} 
[(-1)^{u \cdot w'}
- (-1)^{v \cdot w'}]\\
&=(A_{Q_n} (E_u-E_v))_w
=(n-2k) (E_u-E_v)_w. \qedhere
\end{align*}
\end{proof}

\end{document}
